\chapter{Polynomial comparison for every global ternary coefficient}
\label{chap:bd26-polynomial-full}

\section{The weighted edge integral}

The family $q_{12}(1-q_{12})/(z_1-z_2)$ forces a nonzero image of an unmerged one-form.
The weighted edge integral supplies that image.
It cancels the regular composite states in the adjacent comparison defect before any division by the level.

Retain the base $S$, the ordered complex $O$, the five-partition complex $(E,Q)$, and the chiral operations of Chapter~\ref{chap:bd26-comparison}.
In particular, the source coefficients are global compatible families.
Each adjacent edge ring is $R[t^{-1}]$ and has no outer pole.
All principal parts and all differential-operator sums remain finite.

On an oriented edge with coordinate $u=q_{ij}$, define
\begin{equation}\label{eq:bd26-weighted-integral}
 M(g)=\int_0^1u g(u)\,du,
 \qquad v_g=\pp_t\bigl(p(t)M(g)\bigr),
 \qquad t=z_i-z_j.
\end{equation}
The degree of $M(g)$ is zero after the coefficient one-form $gdu$ has been integrated.
Every normal coefficient of $v_g$ is polynomial in the two remaining block positions.
Polynomial integration by parts gives
\begin{equation}\label{eq:bd26-weighted-endpoint}
 M(\partial_uf)=f(1)-A(f),\qquad
 v_{\partial_uf}+q_{f/t}=\pp_t\bigl(p(t)f(1)\bigr).
\end{equation}
The second identity holds at every finite inner pole order.
It is an identity before any outer collision is taken.

For $c\in C^1$ and the word $e_{ijk}$, write its two actual edge restrictions as
\[
 c|_{\{ijk,jik\}}=g_Ldu_L,
 \qquad c|_{\{ijk,ikj\}}=g_Rdu_R,
 \qquad u_L=q_{ij},\quad u_R=q_{jk}.
\]
Both edges have $u=1$ at the vertex $ijk$.
Let $v_L$ and $v_R$ be their principal parts from \eqref{eq:bd26-weighted-integral}, in their respective normal coordinates.
Define a degree-zero $R$-linear map $B^+:O\to E$ by retaining every mixed component of $B$ and setting
\begin{equation}\label{eq:bd26-positive-leading}
 \begin{aligned}
 B^+(fe_{ijk})&=f(ijk)[J_i|J_j|J_k] && (f\in C^0),\\
 B^+(ce_{ijk})&=[v_L|J_k]-[J_i|v_R] && (c\in C^1),\\
 B^+(ce_{ijk})&=0 && (|c|\geq2).
 \end{aligned}
\end{equation}
The first mixed block in the second line lies on $\Delta_{ij}$.
The second lies on $\Delta_{jk}$.
Each has total degree $-2$, equal to the degree of $ce_{ijk}$.
The minus sign in the right term is the sign of the adjacent $jk$ collision when the remaining suspended current is placed first.

\section{Divisibility of the complete comparison defect}

The defect of $B^+$ is
\[
 D^+=QB^+-B^+b.
\]
It satisfies $QD^++D^+b=0$.
For an outer operation, let a finite inner principal part be called regular in the outer variable when its normal coefficients are polynomials in the block positions.
This property concerns its coefficients before the outer chiral operation.
The outer output itself can have poles.

\begin{lemma}\label{lem:bd26-arbitrary-regular-outer}
For any finite inner principal part $v$ with states in $V$ that is regular in the outer variable,
\[
 R_L(v)\in\kappa P_3(V),\qquad R_R(v)\in\kappa P_3(V).
\]
\end{lemma}
\begin{proof}
Reduce its monomial-state and normal coefficients modulo $\kappa$.
The outer fields are $Y_0(a,s)b=(e^{sT}a)b$.
The coefficient is polynomial in $s$, so the outer principal part is zero.
Every retained inner normal derivative preserves this conclusion, since it differentiates regular coefficients or remains a normal derivative in the final transfer module.
The output therefore vanishes modulo $\kappa$.
The free basis \eqref{eq:bd26-small} gives the asserted divisibility.
\end{proof}

\begin{proposition}\label{prop:bd26-complete-defect}
Every value of $D^+$ belongs to $\kappa E$.
Its values on unmerged words are
\begin{equation}\label{eq:bd26-positive-defect}
 \begin{aligned}
 D^+(fe_{ijk})&=-[\mu_{ik}(f(ijk)(J_i,J_k))|J_j],
                    &&f\in C^0,\\
 D^+(ce_{ijk})&=[R_L(v_L)-R_R(v_R)],
                    &&c\in C^1,\\
 D^+(ce_{ijk})&=0, &&|c|\geq2.
 \end{aligned}
\end{equation}
In the first line, the coefficient $f(ijk)$ is retained as a function of all three positions during the first collision.
The first output lies on $\Delta_{ik}$.
The second output lies on the small diagonal, with both terms written using the same density and normal-coordinate identification.
\end{proposition}
\begin{proof}
For a zero-form $f$, the $ij$ component of $QB^+(fe_{ijk})$ is
$[\pp_t(p(t)f(ijk))|J_k]$.
The corresponding component of $B^+b(fe_{ijk})$ is
\[
 [v_{\partial_uf}+q_{f/t}|J_k].
\]
The value of $f$ at $u=1$ on this edge is $f(ijk)$.
Equation~\eqref{eq:bd26-weighted-endpoint} makes the two components equal.

For the $jk$ component, both expressions have a minus sign with $J_i$ in the first block.
The same endpoint identity proves their equality.
The source differential has no $ik$ face, so the nonadjacent chiral term remains.
The three-current unshuffle sign for that term is negative.
This proves the first line of \eqref{eq:bd26-positive-defect}.

The coefficient $f(ijk)$ is regular in every normal variable.
Thus \eqref{eq:bd26-field} gives
\[
 \mu_{ik}(f(ijk)(J_i,J_k))
 =\kappa\pp_{t_{ik}}\frac{f(ijk)}{t_{ik}^2}\mathbf1.
\]
This proves divisibility on every unmerged zero-form.

For $c\in C^1$, the states in $B^+(ce_{ijk})$ are even and have zero internal differential.
Their collision differential is $[R_L(v_L)-R_R(v_R)]$.
The coefficient differential $d_Cc$ has degree two, so its unmerged image is zero.
The other boundary coefficients are $c\omega$, whose degree-two edge restrictions are zero.
Hence $B^+b(ce_{ijk})=0$.
This proves the second line of \eqref{eq:bd26-positive-defect}.
Lemma~\ref{lem:bd26-arbitrary-regular-outer} proves divisibility of both terms.
All higher-degree unmerged inputs have zero defect by the same edge-degree argument.

The mixed values of $D^+$ equal those of $D$.
On mixed zero-forms, their possible outer term contains $r_f$, a vacuum principal part with polynomial outer coefficients.
The vacuum operation with that outer coefficient has no pole, so this term is zero.
On mixed one-forms, the values are $-[R_L(q_g)]$ and $+[R_R(q_g)]$.
They are divisible by $\kappa$ by Lemma~\ref{lem:bd26-global-divisibility}.
Higher-degree mixed restrictions vanish.
This checks every summand of $O$.
\end{proof}

\section{The polynomial ternary map}

Every summand of $E$ is torsion-free over $S$.
Proposition~\ref{prop:bd26-complete-defect} therefore defines a unique degree-one $R$-linear map
\[
 T^+:O\longrightarrow E,
 \qquad \kappa T^+=D^+.
\]
Its values are obtained by division of polynomial coefficients known to be divisible by $\kappa$.
This definition uses no inverse of $\kappa$ in the coefficient ring.

\begin{theorem}\label{thm:bd26-polynomial-complete}
The formula
\begin{equation}\label{eq:bd26-polynomial-complete}
 G=B^++\mathsf eT^+:O\longrightarrow E
\end{equation}
defines a degree-zero $R$-linear cochain map over $S$.
It is defined on all six unmerged words and all twelve mixed words, with every global compatible-family coefficient and every finite inner pole order.
Its mixed components are exactly those of \eqref{eq:bd26-leading-map} and \eqref{eq:bd26-mixed-correction}, with the polynomial quotients taken before specialization.
Its separated projection is $fe_{ijk}\mapsto f(ijk)[J_i|J_j|J_k]$.
Its positive-degree unmerged components are determined by \eqref{eq:bd26-positive-leading} and \eqref{eq:bd26-positive-defect}.
It specializes to every numerical level, including zero.
\end{theorem}
\begin{proof}
The closed-defect equation and $S$-torsion-freeness imply
\[
 QT^++T^+b=0.
\]
Using $Q\mathsf e+\mathsf eQ=-\kappa$ gives
\[
 \begin{aligned}
 QG-Gb
 &=D^++Q\mathsf eT^+-\mathsf eT^+b\\
 &=\kappa T^+-\kappa T^+-\mathsf e(QT^++T^+b)=0.
 \end{aligned}
\]
This proves the chain equation over $S$ directly.
The complete list of defect components proves the support and component assertions.
Every scalar quotient, integration, and residue is $R$-linear.
All outputs have finite normal expansions.
Finally, specialize the polynomial identity $QG=Gb$ along $S\to S/(\kappa-k)$.
No division at the numerical value $k=0$ occurs.
\end{proof}

In particular, the unmerged one-form image is
\begin{equation}\label{eq:bd26-polynomial-oneform}
 G(ce_{ijk})=[v_L|J_k]-[J_i|v_R]
   +\left[\epsilon\,\frac{R_L(v_L)-R_R(v_R)}\kappa\right].
\end{equation}
The displayed quotient is the unique polynomial quotient.
For a constant unmerged word, $G$ has the same nonadjacent correction as \eqref{eq:bd26-constant-word}.

The necessity of the positive-degree component has a direct check.
Take $f=q_{12}(1-q_{12})/t$ and $e=e_{123}$.
Then $G(fe)=0$ because every vertex value vanishes.
On the left edge,
\[
 M(\partial_uf)=-\frac1{6t},\qquad
 v_{\partial_uf}=-\frac\kappa{6t^3}\mathbf1-\frac{J^2}{6t}
 =-q_{f/t}.
\]
The right edge restriction is zero.
Therefore the proper-diagonal terms of
$G((d_Cf)e)-G(f\omega_{12}m_{12|3})$ cancel.
Their small-diagonal terms also cancel, by the same identity after the outer operation.
Thus $G(b(fe))=0=QG(fe)$, including the nonzero composite state modulo $\kappa$.

\begin{corollary}\label{cor:bd26-polynomial-equivariant}
Replacing $\mathsf e$ in \eqref{eq:bd26-polynomial-complete} by
$(\mathsf e_1+\mathsf e_2+\mathsf e_3)/3$ gives a permutation-equivariant polynomial cochain map $G_{\mathrm{sym}}$.
There are also polynomial cochain maps of right differential modules
\[
 \operatorname{Ind}(O)\longrightarrow E,
 \qquad m\otimes P\longmapsto G(m)P
 \quad\hbox{or}\quad G_{\mathrm{sym}}(m)P.
\]
\end{corollary}
\begin{proof}
Each scalar insertion has anticommutator $-\kappa$ with $Q$.
Their average has the same identity and commutes with label permutations.
The two oriented edge integrals in $B^+$ commute with relabeling, since each has its ordered-word vertex at $u=1$.
The defect and its unique polynomial quotient are therefore equivariant.
This proves the first assertion.
The induction formulas are balanced over $R$ and commute with the differential by Proposition~\ref{prop:bd26-induced-comparison}.
\end{proof}

\section{A comparison with the unextended level-zero chiral algebra}

Put $V_0=\mathbb C[x_1,x_2,\ldots]$ with field
$Y_0(a,t)b=(e^{tT}a)b$ and zero state differential.
Its chiral operation on pole-valued inputs is still the principal-part operation \eqref{eq:bd26-chiral}.
Let $E_0(V_0)$ denote its five-partition Chevalley--Cousin complex, using the same shifts and densities.
Let $O_0=O\otimes_S S/(\kappa)$.

The projection $W\otimes_S S/(\kappa)\to V_0$ sends $\epsilon$ to zero and preserves every vertex operation.
Applying it to each partition factor gives a cochain map to $E_0(V_0)$.
Its composite with the specialization of $G$ is denoted $G_0$.

\begin{corollary}\label{cor:bd26-native-zero}
There is a permutation-equivariant cochain map
\[
 G_0:O_0\longrightarrow E_0(V_0)
\]
to the unextended level-zero chiral algebra.
It is obtained from $B^+$ by setting $\kappa=\epsilon=0$.
It also induces a cochain map of right differential modules from $\operatorname{Ind}(O_0)$.
The degree-$-3$ source cycle
\begin{equation}\label{eq:bd26-nonzero-cycle}
 z=(z_1-z_2)(z_2-z_3)e_{123}
   +(z_2-z_3)q_{12}m_{12|3}
   +(z_1-z_2)q_{23}n_{1|23}
\end{equation}
has image representing a nonzero class in $H^{-3}(E_0(V_0))$.
\end{corollary}
\begin{proof}
The state projection is a differential graded vertex map because both differentials are zero and the scalar factor $\epsilon$ is an ideal.
It therefore commutes with the chiral operation, its transfers, and the Chevalley--Cousin differential.
All correction terms in $G-B^+$ contain $\epsilon$ and vanish under the projection.
Hence $G_0$ is exactly the displayed specialization of $B^+$.
The same specialization of the equivariant map $G_{\mathrm{sym}}$ gives $G_0$, proving equivariance.
Differential induction applies as before.

Equation~\eqref{eq:bd26-source} and $d_Cq_{ij}=(z_i-z_j)\omega_{ij}$ give $bz=0$.
The mixed zero-form images contain $\epsilon$ and therefore vanish under the projection.
Consequently,
\[
 G_0(z)=(z_1-z_2)(z_2-z_3)[J_1|J_2|J_3].
\]
This is nonzero in the separated summand.
It is closed because every regular-input level-zero chiral bracket vanishes.
The complex $E_0(V_0)$ has terms only in degrees $-3,-2,-1$.
There is no incoming boundary in degree $-3$.
Thus the displayed image represents a nonzero cohomology class.
\end{proof}

This comparison uses the unextended state space only at level zero.
It asserts no quasi-isomorphism of the full complexes.

The source of this polynomial map is the global ordered coefficient complex.
Admitting independent outer poles changes the comparison problem.
The cubic obstruction \eqref{eq:bd26-outer-obstruction} still excludes the specified extension on that larger mixed domain with its intermediate images fixed.
The polynomial family before the level-zero projection uses the Koszul state extension $W$.
It preserves neither the nonzero-level diagonal state cohomology of $V$ nor a presently unspecified system of state-translation relations on $O$.
Compatibility with ordered surjections, deconcatenation, and four-input operations remains an additional requirement.
